This paper was introduced to separate the elementary visual and projective
prelude from the heavier foundational machinery of Paper~I.  The division of
labor is deliberately conservative: Paper~I's marked-history terminology and
proved structural results are not rewritten here.

\subsection{What Paper 0 establishes}

The proved elementary spine is
\[
\begin{gathered}
  \text{left and right comb grammars}
  \longrightarrow
  \text{slot-$(1)$/slot-$(2)$ one-hole maps},\\
  \text{affine point path}
  \quad\text{and}\quad
  \text{projective ripple pencil},\\
  \text{reciprocal}
  \longrightarrow
  \text{pole and infinity}
  \longrightarrow
  \text{fixed-point recursion},\\
  \text{one-hole word}
  \longrightarrow
  \text{$2\times2$ prefix matrix}
  \longrightarrow
  \text{path, pole, fixed points, and ripples},\\
  \Aff(1,K)=\Stab(\infty)
  \subset
  \PGL_2(K).
\end{gathered}
\]
The paper also supplies two explicit laboratories:

\begin{itemize}
  \item the zeroth-kind upper-half-plane model, used only to make affine point
  propagation visible;
  \item the pulled-back horocycle pencil, used to make finite poles and
  right-slot reciprocal structure visible.
\end{itemize}

The finite geometric-series and continued-fraction examples are exact at every
truncation.  Their infinite values are asserted only under the stated formal or
analytic convergence hypotheses.

\subsection{What remains in Paper I}

Paper~I begins one level later.  It retains responsibility for:

\begin{enumerate}
  \item the intrinsic dependency-poset classification of all sequential binary
  trees, including mixed slot words;

  \item marked planar sequential trees and bounded/free marked spinal
  histories;

  \item the distinction among planar mirror, operand-slot exchange, temporal
  reversal, and path inverse;

  \item projective evaluation of chronological marked histories, including
  ordinary-domain data and the $q=4$ Hecke sublanguage;

  \item affine cocycles, continuous arithmetic flow, and the intrinsic
  definition of a regular arithmetic expression space;

  \item the complete homogeneous hyperbolic model, regular-zero rigidity,
  global ACS torsion, and contact curvature.
\end{enumerate}

Thus Paper~0 does not replace the formal history layer.  It explains why two
pure comb orientations lead first to path and ripple pictures and why matrices
are the correct common language before mixed chirality is introduced.

\subsection{The revised dependency architecture}

The intended series dependency is now
\[
  \boxed{
  \text{Paper 0}
  \longrightarrow
  \text{Paper I}
  \longrightarrow
  \{\text{Papers II, III, IV}\}.
  }
\]
The first arrow exports the following interfaces:

\begin{center}
\renewcommand{\arraystretch}{1.3}
\begin{tabular}{ll}
\toprule
Paper 0 object & Paper I use \\
\midrule
left/right comb convention & operand-slot chirality on sequential spines \\
point path & bounded marked-history evaluation in the affine sector \\
ripple pencil and pole & motivation for bilateral projective continuation \\
reciprocal and infinity & ordinary/projective domain distinction \\
fixed-point recursion & first examples separating closure from identity \\
prefix matrices & chronological projective evaluation \\
$\Aff(1)=\Stab(\infty)$ & affine flow as the Borel sector of projective AEG \\
\bottomrule
\end{tabular}
\end{center}

Paper~I may repeat a concise definition or theorem when needed for logical
self-containment, but Paper~0 is the canonical source for the elementary
path/ripple exposition and the reciprocal--pole--infinity discussion.  In
particular, later revisions of Paper~I should avoid maintaining a second,
divergent pedagogical treatment of those themes.

\subsection{Open extensions}

The elementary synthesis leaves several genuine problems.

\begin{enumerate}
  \item Construct a projective analogue of the affine arithmetic expression
  space in which the pole and ripple pencil are intrinsic rather than chosen
  through one upper-half-plane chart.

  \item Determine which projective path/ripple data descend through relations
  among history words and which require a lift to a full projective frame.

  \item Extend the one-hole analysis to multi-wire expression networks while
  retaining ordinary-domain, pole, and fixed-point information.

  \item Develop a controlled theory of infinite arithmetic expressions that
  separates formal convergence, analytic convergence, projective convergence,
  and convergence of realized computation.
\end{enumerate}

None of these problems is solved merely by the matrix representation.  The
stable conclusion of the present paper is more elementary: the straight path
at a fixed pole and the nested ripple pencil at a moving finite pole are two
faces of one projective arithmetic word.
